\subsection{Числовые ряды. Абсолютная и условная сходимость. Критерий Коши. Признак сравнения.}

\subsubsection*{Определения и элементарные факты}
Пусть $\{a_k\}$ — числовая последовательность.
\[ S_n = \sum^n_{k=1}a_k \]
\begin{tcolorbox}[title=Определения]
    \textbf{Ряд} — это символ $\sum^\infty_{k=1} a_k$. \\
    $a_k$ — \textbf{общий член ряда}.\\
    $S_n$ — \textbf{n-я частичная сумма ряда}.\\
    
    \begin{itemize}
        \item Если существует конечный предел $S = \lim_{n\to\infty} S_n$, то ряд называется \textbf{сходящимся}, а $S$ — его \textbf{суммой} ($\sum_{k=1}^\infty a_k = S$).
        \item Если предел $S_n$ не существует или бесконечен, то ряд называется \textbf{расходящимся}.
    \end{itemize}
\end{tcolorbox}

\subsubsection*{Теорема 1 (Линейность)}
Если ряды сходятся, то:
\[ \sum_{k=1}^\infty (a_k \pm b_k) = \sum_{k=1}^\infty a_k \pm \sum_{k=1}^\infty b_k \]
\[ \sum_{k=1}^\infty (c \cdot a_k) = c \cdot \sum_{k=1}^\infty a_k \]

\subsubsection*{Теорема 2. Критерий Коши сходимости ряда}
Ряд $\sum a_k$ сходится $\Leftrightarrow \forall \varepsilon > 0\;\; \exists n_\varepsilon : \forall n > n_\varepsilon, \forall p \in \mathbb{N}$:
\[ |S_{n+p} - S_n| = \left| \sum_{k=n+1}^{n+p} a_k \right| < \varepsilon \]

\textbf{Следствие 1.} Изменение конечного числа членов ряда не влияет на его сходимость (но влияет на сумму).

\textbf{Следствие 2 (Необходимое условие сходимости).}
Если $\sum a_k$ сходится $\Rightarrow \lim_{k\to\infty} a_k = 0$.
\begin{tcolorbox}[title=Доказательство]
    Пусть ряд сходится. По критерию Коши для любого $\varepsilon$ (при $p=1$):
    \[ |S_{n+1} - S_n| = |a_{n+1}| < \varepsilon \]
    Это и означает, что $a_n \to 0$.
\end{tcolorbox}

\subsubsection*{Примеры рядов}
\textbf{Пример 1 (Геометрическая прогрессия).}
\[ \sum_{k=0}^\infty q^k, \quad (|q| < 1) \]
\[ S_n = 1 + q + \dots + q^n = \frac{1-q^{n+1}}{1-q} \xrightarrow[n\to\infty]{} \frac{1}{1-q} \]

\textbf{Пример 2 (Гармонический ряд).}
\[ \sum_{n=1}^\infty \frac{1}{n} \quad \text{— расходится.} \]
Доказательство через критерий Коши. Оценим отрезок ряда (от $n+1$ до $2n$, то есть $p=n$):
\[ \left| \sum_{k=n+1}^{2n} \frac{1}{k} \right| = \frac{1}{n+1} + \dots + \frac{1}{2n} > \underbrace{\frac{1}{2n} + \dots + \frac{1}{2n}}_{n \text{ слагаемых}} = n \cdot \frac{1}{2n} = \frac{1}{2} \]
Так как разность частичных сумм не стремится к нулю (всегда $> 1/2$), ряд расходится.

\textbf{Пример 3 (Знакопеременный ряд и перестановки).}
\[ 1 - 1 + \frac{1}{2} - \frac{1}{2} + \frac{1}{3} - \frac{1}{3} + \dots \]
\begin{itemize}
    \item Исходный ряд сходится (сумма 0), т.к. $S_{2n} = 0, S_{2n+1} \to 0$.
    \item \textbf{Перестановка (Теорема Римана):} Если мы переставим члены условно сходящегося ряда, можно получить любую сумму или расходимость.
    Например, наберем столько положительных членов, чтобы сумма превысила 100, потом вычтем один отрицательный, потом снова наберем положительных... Так можно уйти в $+\infty$.
\end{itemize}

\subsubsection*{Абсолютная и условная сходимость}
\begin{tcolorbox}[title=Определения]
    \begin{itemize}
        \item Ряд $\sum a_k$ сходится \textbf{абсолютно}, если сходится ряд из модулей $\sum |a_k|$.
        \item Ряд $\sum a_k$ сходится \textbf{условно}, если сам он сходится, а ряд модулей $\sum |a_k|$ расходится.
    \end{itemize}
\end{tcolorbox}

\textbf{Теорема.} Если ряд сходится абсолютно, то он сходится (обычно).
\begin{tcolorbox}[title=Доказательство]
    По критерию Коши для сходящегося $\sum |a_k|$:
    \[ \forall \varepsilon > 0 \dots \sum_{k=n+1}^{n+p} |a_k| < \varepsilon \]
    Используем неравенство треугольника для суммы:
    \[ \left| \sum_{k=n+1}^{n+p} a_k \right| \leq \sum_{k=n+1}^{n+p} |a_k| < \varepsilon \]
    Значит, критерий Коши выполнен и для исходного ряда $\sum a_k$.
\end{tcolorbox}

\subsubsection*{Ряды с неотрицательными членами}
\textbf{Теорема.} Пусть $a_k \ge 0$. Ряд $\sum a_k$ сходится $\Leftrightarrow$ последовательность частичных сумм $\{S_n\}$ ограничена сверху.
\begin{tcolorbox}[title=Доказательство]
    Так как $a_k \ge 0$, то $S_{n+1} = S_n + a_{n+1} \ge S_n$.
    Последовательность $S_n$ монотонно возрастает.
    По теореме Вейерштрасса: монотонная последовательность сходится $\Leftrightarrow$ она ограничена.
\end{tcolorbox}

\subsubsection*{Признак сравнения}
Пусть даны два ряда $\sum a_k$ и $\sum b_k$ с неотрицательными членами, причем:
\[ 0 \le a_k \le b_k \quad (\forall k) \]
($a_k$ — «меньший» ряд, $b_k$ — «больший» ряд).

\textbf{Теорема.}
\begin{enumerate}
    \item Если сходится **больший** ряд ($\sum b_k$), то сходится и **меньший** ($\sum a_k$).
    \item Если расходится **меньший** ряд ($\sum a_k$), то расходится и **больший** ($\sum b_k$).
\end{enumerate}

\begin{tcolorbox}[title=Доказательство]
    Обозначим частичные суммы $A_n = \sum a_k$ и $B_n = \sum b_k$.
    Так как $a_k \le b_k$, то $A_n \le B_n$.
    \begin{enumerate}
        \item Если $\sum b_k$ сходится, то $B_n$ ограничена. Тогда и $A_n$ ограничена (сверху числом $\lim B_n$). Так как $A_n \uparrow$, то $\sum a_k$ сходится.
        \item Если $\sum a_k$ расходится, то $A_n \to +\infty$. Так как $B_n \ge A_n$, то $B_n \to +\infty$.
    \end{enumerate}
\end{tcolorbox}
 
\textbf{Следствие (для знакопеременных рядов):}
Если $|a_k| \le b_k$ и ряд $\sum b_k$ сходится, то ряд $\sum a_k$ сходится абсолютно.

\textbf{Пример.} $\sum_{n=1}^\infty \frac{\sin n}{n^2}$.
\[ \left| \frac{\sin n}{n^2} \right| \le \frac{1}{n^2} \]
Сравним с рядом $\sum \frac{1}{n^2}$. Он сходится (так как $n^2 > n(n-1)$):
\[ \frac{1}{n^2} < \frac{1}{n(n-1)} = \frac{1}{n-1} - \frac{1}{n} \]
Сумма мажорирующего ряда (Телескопическая сумма):
\[ \sum_{n=2}^\infty \left(\frac{1}{n-1} - \frac{1}{n}\right) = \left(1 - \frac{1}{2}\right) + \left(\frac{1}{2} - \frac{1}{3}\right) + \dots = 1 \]
Ряд $\sum \frac{1}{n(n-1)}$ сходится $\Rightarrow \sum \frac{1}{n^2}$ сходится $\Rightarrow \sum \frac{\sin n}{n^2}$ сходится абсолютно.
